\documentclass[11pt]{article}
\usepackage[T1]{fontenc}
\usepackage[utf8]{inputenc}
\usepackage[margin=1in]{geometry}
\usepackage{amsmath,amssymb,amsthm}
\usepackage{microtype}
\usepackage[hidelinks]{hyperref}
\newcommand{\E}{\mathbb E}
\newcommand{\Pb}{\mathbb P}
\newcommand{\Ind}[1]{\mathbf 1\{#1\}}
\newtheorem{proposition}{Proposition}
\newtheorem{lemma}{Lemma}
\title{Fluid model forgetting case}
\author{Standalone theoretical note}
\date{}
\begin{document}
\maketitle

This note presents the fluid model and some results. 

\subsection*{Fluid bellman for fixed policy U}
Recall $U^1_\rho$ corresponds to the value function when applying the fixed policy of using always product 1 (for simplicity of this note I'm going to use just $U$). For the purpose of this note we are going to work with the normalized coordinates $x = (1-\rho)S$. We start from the bellman equation

\begin{equation}
  U(x) = (1-D(x)) \gamma U(x)+ D(x)[ R-c_1 + \gamma(p_1U(T^s_\rho(x)) + (1-p_1)U(T^f_\rho(x)))]
\end{equation}
Considering the transitions as 
\begin{align}
  T^s_{dt}(x) &= \rho^{dt}x + (1-\rho^{dt}) \\
  T^f_{dt}(x) &= \rho^{dt}x
\end{align}
We set 
\begin{align*}
  \rho^{dt} = e^{-\lambda dt} = 1- \lambda dt + o(dt)\\
  \gamma^{dt} = e^{-\delta dt} = 1- \delta dt + o(dt)
\end{align*}
Then, replacing in the transition functions we obtain
\begin{align}
  T^s_{dt}(x) - x &= \lambda(1- x)dt + o(dt)\\
  T^f_{dt}(x) - x &= -\lambda xdt + o(dt)
\end{align}
These are our ODE equations
\begin{equation}
  \begin{cases}
    \text{success: } &\dot{x} = \lambda(1- x)\\
    \text{failure: } &\dot{x} = -\lambda x 
  \end{cases}
\end{equation}
Now, we note that the expected change of the state in the bellman equation would be 
\begin{equation}
  D(x)[p_1 \lambda(1-x) + (1-p_1)(-\lambda x)] dt = D(x) \lambda (p_1-x)dt
\end{equation}
Hence we have 
\begin{equation}
  \dot{x}= D(x)\lambda (p_1-x) 
\end{equation}
We can write our bellman equation over $dt$ as follows
\begin{equation}
  U(x) = (1-D(x)) \gamma^{dt} U(x)+ D(x)[ (R-c_1)dt + \gamma^{dt}(p_1U(T^s_{dt}(x)) + (1-p_1)U(T^f_{dt}(x)))]
\end{equation}
Now, we approximate $U(T^s_{dt}(x))$ and $U(T^f_{dt}(x))$ using Taylor expansion
\begin{align}
  U(T^s_{dt}(x)) &= U(x) + \lambda(1-x)U'(x)dt + o(dt) \\
  U(T^f_{dt}(x)) &= U(x) - \lambda x U'(x)dt + o(dt) 
\end{align}
replacing the discount factor's approximations 
\begin{align}
  U(x) = (1-D(x))(1-\delta dt)U(x) + D(x)[(R-c_1)dt + \\(1-\delta dt)(U(x) + \lambda(p_1-x)U'(x)dt) ]+ o(dt)
\end{align}
Rearranging terms we get
\begin{align}
  U(x) &= U(x) + [D(x)(R-c_1) + D(x) (\lambda (p_1-x)U'(x))-\delta U(x)]dt  + o(dt)\\ 
  \Rightarrow 0 &= [D(x)(R-c_1) + D(x) (\lambda (p_1-x)U'(x))-\delta U(x)]dt  + o(dt)
\end{align}
Dividing by $dt$ and taking $dt \rightarrow 0$ we get
\begin{equation}
  \delta U(x) = D(x)(R-c_1) + D(x)\lambda (p_1-x)U'(x)
\end{equation}
% Evaluating this at $p_1$ we get 
% \begin{equation}
%   U(p_1) = \frac{D(p_1)(R-c_1)}{\delta} 
% \end{equation}
To solve it, let $x(t)$ denote the fluid state at time $t$, starting from $x(0)=x$. The fluid Bellman equation can therefore be written as
\begin{equation}
    \delta U(x(t))=D(x(t))(R-c_1)+U'(x(t))\dot{x}(t).
\end{equation}
By the chain rule,
\begin{equation}
    \frac{d}{dt}U(x(t))=U'(x(t))\dot{x}(t).
\end{equation}
Hence,
\begin{equation}
    \frac{d}{dt}U(x(t)) - \delta U(x(t)) = -D(x(t))(R-c_1) 
\end{equation}
Multiplying both sides by $e^{-\delta t}$,
\begin{equation}
    e^{-\delta t}\frac{d}{dt}U(x(t))-\delta e^{-\delta t}U(x(t))=-(R-c_1)e^{-\delta t}D(x(t)).
\end{equation}
The left-hand side is the derivative of $e^{-\delta t}U(x(t))$, so
\begin{equation}
    \frac{d}{dt}
    \left[
        e^{-\delta t}U(x(t))
    \right]=-(R-c_1)e^{-\delta t}D(x(t)).
\end{equation}
Integrating from $0$ to $T$,
\begin{equation}
    e^{-\delta T}U(x(T))-U(x)=-(R-c_1)\int_0^T e^{-\delta t}D(x(t))\,dt.
\end{equation}
Rearranging,
\begin{equation}
    U(x)=e^{-\delta T}U(x(T))+(R-c_1)\int_0^T e^{-\delta t}D(x(t))\,dt.
\end{equation}
Since $D(x)>0$, the sign of $\dot{x}(t)$ is determined by $p_1-x(t)$.
Thus, $x(t)$ moves toward $p_1$ from either side and
\begin{equation}
    x(t) \rightarrow p_1
    \qquad \text{as } t\rightarrow\infty.
\end{equation}
Moreover, $U$ is bounded and $\delta>0$, so
\begin{equation}
    e^{-\delta T}U(x(T))\rightarrow 0.
\end{equation}
Taking $T\rightarrow\infty$, we obtain the fluid value under the fixed
policy of always using product 1:
\begin{equation}
    U(x) = (R-c_1) \int_0^\infty e^{-\delta t}D(x(t))dt
\end{equation}
where $x(t)$ solves
\begin{equation}
    \dot{x}(t)=D(x(t))\lambda(p_1-x(t)),\qquad x(0)=x.
\end{equation}

This expression provides an integral representation of the fluid value function. The trajectory $x(t)$ is determined by the demand function
$D(\cdot)$. 

\subsection*{Fluid bellman for V}
We start from the bellman equation 
\begin{equation*}
  V(x)= (1-D(x))\gamma V(x) + D(x)\max_{i \in\{1,2\}} \left\{R-c_i + \gamma[p_iV(T^s(x)) +(1-p_i)V(T^f(x))]\right\}
\end{equation*}
similar as before we can write
\begin{equation*}
  \begin{aligned}
  V(x)={}& (1-D(x))(1-\delta dt) V(x) \\
  &+ D(x)\max_{i \in\{1,2\}} \left\{(R-c_i)dt + (1-\delta dt)[V(x) +\lambda(p_i-x)V'(x)dt]\right\} + o(dt)
  \end{aligned}
\end{equation*}
distributing and rearranging terms we get 
\begin{align*}
  V(x)= (1-D(x))&(1-\delta dt) V(x) + \\D(x)\max_{i \in\{1,2\}} &\left\{(R-c_i)dt + V(x) + [\lambda(p_i-x)V'(x) - \delta V(x)]dt \right\} +o(dt)
\end{align*}
Since $V(x)$ and $\delta V(x)$ do not depend on $i$ we get
\begin{align*}
  V(x)= V(x) + [D(x)\max_{i \in\{1,2\}} &\left\{R-c_i + \lambda(p_i-x)V'(x) \right\} -\delta V(x)]dt + o(dt)
\end{align*}
dividing by $dt$ and taking $dt\rightarrow 0 $ we get
\begin{align*}
  \delta V(x)= D(x)\max_{i \in\{1,2\}} &\left\{R-c_i + \lambda(p_i-x)V'(x) \right\} 
\end{align*}
We get that product 2 is optimal when 
\begin{eqnarray}
  R-c_2 + \lambda(p_2-x)V'(x) &>& R-c_1 + \lambda(p_1-x)V'(x)\\
  \iff V'(x) &>& \frac{\Delta c}{\lambda\Delta p}
\end{eqnarray}
Which is basically the same equation we had in the original model. 


\end{document}
